% Figure: thin cylindrical pressure vessel showing hoop stress (around the
% circumference) and axial stress (along the axis), with the 2:1 ratio.
\begin{figure}[htbp]
\centering
\begin{tikzpicture}[scale=1.0,>=latex]
  % cylinder body
  \draw[engblack, thick] (0,-1) -- (5,-1);
  \draw[engblack, thick] (0,1) -- (5,1);
  % left end ellipse
  \draw[engblack, thick] (0,0) ellipse [x radius=0.45, y radius=1];
  % right end ellipse (front half solid)
  \draw[engblack, thick] (5,0) ellipse [x radius=0.45, y radius=1];
  % hoop stress arrows around a ring (drawn as a band at x=2.5)
  \draw[engblue, very thick, dashed] (2.5,0) ellipse [x radius=0.4, y radius=1];
  \draw[engblue, ultra thick, ->] (2.5,1) -- (2.9,1.0);
  \draw[engblue, ultra thick, ->] (2.5,-1) -- (2.1,-1.0);
  \node[engblue, font=\small] at (2.5,1.4) {$\sigma_\theta = pR/h$ (hoop)};
  % axial stress arrows along the axis
  \draw[engred, ultra thick, ->] (5.0,0) -- (5.9,0);
  \draw[engred, ultra thick, ->] (0.0,0) -- (-0.9,0);
  \node[engred, font=\small] at (2.5,0.0) {$\sigma_z = pR/2h$ (axial)};
  % internal pressure indicator
  \foreach \a in {30,90,150,210,270,330}{
    \draw[englime,->] (2.5,0) -- ++({0.55*cos(\a)},{0.55*sin(\a)});
  }
  \node[englime, font=\scriptsize] at (3.6,-1.35) {internal pressure $p$};
  % radius callout
  \draw[enggray,<->] (4.55,0) -- (4.55,1) node[midway,right,font=\scriptsize]{$R$};
\end{tikzpicture}
\caption[Hoop and axial stress in a cylindrical vessel]{Membrane
stresses in a thin cylindrical pressure vessel of radius $R$, wall
thickness $h$, internal gauge pressure $p$. The circumferential
(hoop) stress $\sigma_\theta = pR/h$ acts around the ring; the axial
stress $\sigma_z = pR/2h$ acts along the axis. The hoop stress is
exactly twice the axial stress, which is why a thin cylinder under
internal pressure splits along a longitudinal seam rather than
parting at a circumferential one, and why the longitudinal weld is
the governing detail in vessel design.}
\label{fig:vol03ch09:cylinder-stresses}
\end{figure}
